\section{Boundary applications beyond grammar}\label{sec:3}

The schema can also describe form--meaning access when $R$ is interpretation, learning, or processing rather than agreement. Constructional compositionality \autocite{goldberg-2006-constructions-at-work}, rated transparency in noun--noun compounds \autocite{bell-schafer-2016-modelling-transparency}, and masked morphological priming \autocite{heyer-kornishova-2018-priming-eventually} supply possible applications. In each case, the analyst would need to identify a bearer, a distinct mediator, a relation, and a constructional counterpart that could conceal the relevant property.

These examples mark scope rather than deliver additional worked profiles. Their relations, items, diagnostics, and proposed stabilizers differ, and the available studies don't establish a common within-items cluster. Constructional interpretation, compound ratings, and morphological priming belong to a research programme outside the evidential spine of this paper. The examples show why the schema keeps $R$ explicit: a construction may be accessible for one relation and opaque for another. They don't warrant a cross-domain projectibility claim here.
